\section{Excitation covariates: the Volterra system that arises}
\label{sec:excitation-cov}

We now let covariates modulate the \emph{excitation itself} --- the strength with
which one event triggers others. This is the genuinely harder case, and the whole
point of developing the Volterra viewpoint: the governing equation stops being a
convolution.

\subsection{The model}

\begin{model}[Covariate-modulated excitation]\label{mod:excitation}
Let observed covariates $Z(t)\in\R^q$ scale the triggering strength through a
log-linear link,
\begin{equation}
\alpha_{mj}(t)=\alpha^{0}_{mj}\,\exp\!\bigl(\delta_{mj}^{\top}Z(t)\bigr),
\qquad
\phi_{mj}(t,u)=\alpha_{mj}(t)\,g_{mj}(t-u),
\label{eq:excit-kernel}
\end{equation}
where $g_{mj}\ge0$ is a fixed shape (e.g.\ $g_{mj}(\tau)=\theta_{mj}e^{-\theta_{mj}\tau}$)
and $\alpha^0_{mj}\ge0$ is the covariate-free base strength.
\end{model}

The modulation may be read at the \emph{receiving} time $t$ (as written) or at the
\emph{parent} time $u$ (replace $Z(t)$ by $Z(u)$); both are handled below and the
implementation uses the former.

\subsection{The mean-intensity equation is now a general Volterra equation}

Averaging the Hawkes intensity exactly as in Proposition~\ref{prop:mbpp-eq}, but
with the time-dependent kernel \eqref{eq:excit-kernel}, gives

\begin{proposition}[Non-convolution MBPP equation]\label{prop:nonconv}
Under (A1)--(A7), the expected intensity solves the system of linear Volterra
equations of the second kind
\begin{equation}
\xi_m(t)=s_m(t)+\sum_{j=1}^M\int_0^t K_{mj}(t,u)\,\xi_j(u)\dd u,
\qquad
K_{mj}(t,u)=\alpha^0_{mj}\,e^{\delta_{mj}^\top Z(t)}\,g_{mj}(t-u).
\label{eq:mbpp-nonconv}
\end{equation}
The kernel $K_{mj}(t,u)$ depends on $t$ and $u$ \emph{separately}, not only on
$t-u$.
\end{proposition}

\begin{intuition}
In \S\ref{sec:no-covariates}--\ref{sec:baseline-cov} the kernel was a function of
the elapsed time $t-u$ alone --- the system reacted to a stimulus the same way
regardless of \emph{when} it arrived. Now the response strength rides on $Z(t)$, so
\emph{when} matters: the same past event echoes more loudly in a high-$Z$ period
than in a low-$Z$ period. Mathematically the kernel is no longer a convolution, and
\textbf{this is exactly what costs us the LTI toolkit}.
\end{intuition}

\subsection{What survives, and what is lost}

\begin{description}[leftmargin=1.4em,style=nextline]
\item[Well-posedness survives.] Equation~\eqref{eq:mbpp-nonconv} is a standard
linear Volterra system. By the general theory (Appendix~\ref{app:volterra}, which
does \emph{not} use convolution structure), for any locally bounded kernel and
continuous $s$ it has a unique solution on every finite horizon, given by the
Neumann series $\xi=\sum_{n\ge0}\mathcal K^n s$ with
$(\mathcal K\xi)_m(t)=\sum_j\int_0^t K_{mj}(t,u)\xi_j(u)\dd u$.
\item[Time-invariance is lost.] Because $K(t,u)\ne K(t-u)$, the convolution theorem
no longer applies: there is \emph{no} Laplace/Fourier solution, \emph{no} constant
transfer function $R(\omega)$, and the spectral (Fourier-neural-operator) shortcut
of the convolution case is no longer exact. The resolvent is now a genuine
two-variable kernel $R(t,u)$ solving $R=K+K\!\circ\!R$ (Appendix~\ref{app:volterra}),
not a single function of $t-u$.
\end{description}

\subsection{Exponential shape $\Rightarrow$ a finite linear time-varying ODE}

All is not lost computationally. The key is that an exponential $g$ makes the
kernel \emph{separable} (degenerate): with $g_{mj}(\tau)=e^{-b_{mj}\tau}$,
\begin{equation*}
K_{mj}(t,u)=\underbrace{\alpha^0_{mj}\,e^{\delta_{mj}^\top Z(t)}\,e^{-b_{mj}t}}_{a_{mj}(t)}\;\underbrace{e^{\,b_{mj}u}}_{\text{(}u\text{ only)}},
\end{equation*}
a rank-one product of a $t$-function and a $u$-function. A separable Volterra
kernel collapses the integral equation onto finitely many ODEs
(Appendix~\ref{app:volterra}). Concretely, introducing the filtered histories
$y_{mj}(t)=\int_0^t e^{-b_{mj}(t-u)}\xi_j(u)\dd u$:

\begin{theorem}[LTV state-space reduction]\label{thm:ltv}
For exponential shapes, the non-convolution MBPP equation~\eqref{eq:mbpp-nonconv}
(receiving-time modulation) is equivalent to the finite \emph{linear time-varying}
ODE system
\begin{equation}
y_{mj}'(t)=\xi_j(t)-b_{mj}\,y_{mj}(t),
\qquad
\xi_m(t)=s_m(t)+\sum_{j} \alpha^0_{mj}\,e^{\delta_{mj}^\top Z(t)}\,y_{mj}(t),
\label{eq:ltv}
\end{equation}
with $y_{mj}(0)=0$. The coefficients are constant except for the modulation
$e^{\delta^\top Z(t)}$, which makes the system time-\emph{varying} rather than
time-invariant.
\end{theorem}

\begin{intuition}
The history of each component is still summarised by a finite ``running, decaying
memory'' $y_{mj}$ --- exactly the state of the convolution case
\eqref{eq:statespace}. The \emph{only} change is that the gain with which that
memory feeds back, $\alpha^0_{mj}e^{\delta^\top Z(t)}$, now breathes in and out with
the covariate. So we keep a finite-dimensional, integrable system; we merely lose
the constant-coefficient (and hence transform-solvable) structure. ``Same state,
time-varying gain.''
\end{intuition}

The system \eqref{eq:ltv} is integrated by Runge--Kutta exactly like the
convolution case; the package provides it as \code{solve\_mbpp\_ltv}. Two checks
confirm correctness: with $Z\equiv 0$ (no modulation) it reproduces the LTI
multivariate solver to machine precision ($\sim10^{-16}$); and with a covariate
that, say, \emph{doubles} the excitation over a window, the intensity rises sharply
during the window and \emph{relaxes back} to the unmodulated stationary level
afterwards --- the signature of a stable forced system.

\begin{remark}[Power-law shapes: direct Volterra quadrature]
If $g$ is \emph{not} a finite sum of exponentials (e.g.\ power-law), the kernel is
\emph{not} separable and no finite ODE exists. One then solves
\eqref{eq:mbpp-nonconv} directly by discretising the Volterra integral on the
observation grid --- a triangular (lower-block-triangular) linear system solved by
forward substitution in $O(N^2 M^2)$ (the Nyström method,
Appendix~\ref{app:volterra}). One may also approximate $g$ by a sum of exponentials
(Prony) to restore a finite LTV ODE to any desired accuracy.
\end{remark}

\subsection{Fitting with excitation covariates}

The MBPP is \emph{still} a Poisson process with a deterministic (now
harder-to-compute) intensity, so the interval-censored likelihood~\eqref{eq:icll}
still applies verbatim: assemble $\Xi(o_{i-1},o_i]$ by integrating the LTV ODE
\eqref{eq:ltv} (or by quadrature), and minimise $\mathcal L_{\mathrm{IC}}$ over
$(\alpha^0,\theta,\delta,\gamma)$. Three things change relative to the convolution
case:
\begin{enumerate}[leftmargin=2.0em]
\item \textbf{More parameters.} Each excitation-covariate coefficient $\delta_{mj}$
adds dimensions; with $p$ covariates and an $M\times M$ network that is up to
$pM^2$ extra parameters. An $\ell_1$ penalty or low-rank/shared-$\delta$ structure
is advisable.
\item \textbf{Non-convexity.} The $\delta$ enter through $e^{\delta^\top Z}$
\emph{inside} the intensity, so the loss is non-convex (unlike the baseline-only
case, which is well-behaved). Use multiple restarts.
\item \textbf{Identification needs interaction.} A $\delta_{mj}$ is identified only
if component $j$ actually fires while $Z$ varies; without that co-occurrence the
excitation modulation leaves no fingerprint in the counts.
\end{enumerate}

\begin{intuition}
We have traded a closed-form, convex, transform-solvable problem for a
finite-dimensional but \emph{time-varying} and \emph{non-convex} one. The physics
is appealing --- ``contagion that is more contagious in some regimes than others''
--- and entirely doable; one simply pays for it in computation (integrate an ODE
instead of evaluating a formula) and in statistics (more parameters, a rougher
landscape, a stronger need for covariate variation). That trade is the precise
content of ``excitation covariates are harder.''
\end{intuition}

\subsection{Summary of the three regimes}

\begin{center}
\renewcommand{\arraystretch}{1.3}
\begin{tabular}{@{}lccc@{}}
\toprule
 & \S\ref{sec:no-covariates}/\ref{sec:baseline-cov} no/baseline cov.\ & \multicolumn{2}{c}{\S\ref{sec:excitation-cov} excitation covariates}\\
\cmidrule(l){3-4}
 & convolution & exp.\ shape & general shape\\
\midrule
MBPP equation & convolution Volterra & \multicolumn{2}{c}{general Volterra (non-convolution)}\\
Time-invariant? & yes (LTI) & no & no\\
Laplace/Fourier solution & yes & no & no\\
Reduces to finite ODE? & yes (constant coef.) & yes (\emph{time-varying} coef.) & no\\
How we compute $\xi,\Xi$ & closed form & integrate LTV ODE & Volterra quadrature\\
Fitting loss & convex(ish) & non-convex & non-convex\\
\bottomrule
\end{tabular}
\end{center}

\noindent
Reading the table left to right is reading this document: each added layer of
realism removes a layer of mathematical convenience, and the Volterra formulation
is what lets us track exactly which convenience is lost and what replaces it.
